% Figura: filtro activo paso bajo de segundo orden (figura 1 del guion), con nudos y corrientes.
\begin{figure}[ht]
\centering
\begin{circuitikz}[circuito]
  \draw (0,0) node[op amp, ao, noinv input up] (A) {};
  % salida
  \draw (A.out) -- ++(0.9,0) coordinate (O);
  \draw (O) to[short, *-o] ++(0.8,0) node[right, tens] {$\Vout$};
  % entrada no inversora: nudo P, C2 a tierra
  \draw (A.+) -- ++(-1.5,0) coordinate (P);
  \draw (P) to[C, l_=$C_2$, i>_=$I_{C2}$] ++(0,-2.4) node[ground] {};
  % R2 entre X y P
  \draw (P) to[R, l_=$R_2$, i<_=$I_2$] ++(-2.2,0) coordinate (X);
  % R1 entre la entrada y X
  \draw (X) to[R, l_=$R_1$, i<^=$I_1$] ++(-2.3,0) coordinate (E);
  \draw (E) to[short, -o] ++(-0.2,0) coordinate (Ein) node[above, tens] {$\Vin$};
  \draw (Ein) -- ++(-0.7,0) to[sV, l_=$v_{\mathrm{in}}$] ++(0,-2.6) node[ground] {};
  % C1 entre X y la salida
  \draw (X) -- ++(0,1.9) coordinate (T) to[C, l=$C_1$, i>^=$I_{C1}$] (T -| O) -- (O);
  % red RA-RB en la entrada inversora
  \draw (A.-) -- ++(-0.4,0) -- ++(0,-1.3) coordinate (N);
  \draw (N) to[R, l_=$R_A$] (N -| O) -- (O);
  \draw (N) to[R, l=$R_B$] ++(0,-1.6) node[ground] {};
  % nudos y tensiones
  \fill (X) circle (1.6pt); \fill (P) circle (1.6pt); \fill (N) circle (1.6pt);
  \node[tens, anchor=north west, font=\scriptsize] at ($(X)+(0.06,-0.06)$) {$V_x$};
  \node[tens, anchor=south west, font=\scriptsize] at ($(P)+(0.06,0.06)$) {$V_+$};
  \node[tens, anchor=north east, font=\scriptsize] at ($(N)+(-0.06,-0.06)$) {$V_-$};
\end{circuitikz}
\caption{Filtro activo paso bajo de segundo orden (figura~1 del guion), con la alimentación omitida.
Se indican los nudos del análisis ($V_x$, $V_+$, $V_-$) y los sentidos de corriente elegidos.}
\label{fig:filtro}
\end{figure}
